\section{The 2026--27 prospective test}
\label{sec:live}

Every season through 2026 has now been used at game level by some part of
\neurhl{}, so the 2026--27 season, which begins on 29 September 2026, is the
only data on which all the models can be tested cleanly. Two sets of
predictions are scored.

\paragraph{Frozen preseason predictions.}
On 25 September 2026, four days before opening night, \neurhl{} 1.0 committed
static predictions with their SHA-256 hashes: a home-win probability for each
of the 1,344 regular-season games, team points with 80\% intervals and
playoff probabilities from 20,000 simulated seasons, and goals, assists and
points for 711 skaters. The scoring plan fixes the statistics in advance (per-game
log loss against the preseason Elo reference; standings error and interval
coverage against a house benchmark; skater points error for the shipped blend
and its two components), states the expected outcomes, including that Elo is
expected to win the per-game comparison because the season engine behind these
probabilities scored worse than Elo on the tune window, and allows one
inference only, after the last regular-season game. After the NHL roster
deadline on 28 September a dated update re-runs the same frozen models on the
final rosters, with availability entries such as a goaltender suspended by his
club removed. It is hashed and scored beside the originals, which remain
primary.

\paragraph{Game-day forecasts.}
\neurhlG{} is live as an exploratory model and \neurhlH{} as the primary win-probability
model. The live engine is the candidate trained through 2024. Its stack,
fitted on out-of-sample predictions for 14,940 games of 2011--2024, puts
coefficients of 0.21, 0.29 and 0.58 on the Elo, \neurhlG{} and \neurhlH{} logits
(0.63 and 0.46 on Elo and \neurhlG{} when \neurhlH{} is unavailable). Each game day produces a morning forecast at about 11:00 Eastern time
with projected lineups and a pregame forecast 45 to 75 minutes before the
scheduled start with the latest lineups and starting goalies. Each forecast
records where its lineup came from: the league's own game data, confirmed
lines, projected lines, or a fallback built from the team's last dressed roster
minus unavailable players. A forecast counts only if the commit that adds it is
on the public repository's main branch before the scheduled start, checked
against the remote and dated independently by a hosted continuous-integration
run. A game without a valid forecast is counted as missed for every model, so
that missing data cannot favour one of them, and nothing is backfilled. The
pregame forecast is primary and the morning forecast secondary. In-season Elo
and \neurhlH{}, fed the same lineups, are scored on the same games, and the frozen
1.0 probabilities are reported descriptively. Inference is made once, after
the last regular-season game, with a bootstrap over weeks of play. With about
1,300 games the standard error of a paired difference is about 0.0024, so only
gaps of about 0.007 or more are likely to be resolved within the season.

\paragraph{The stat sheet's scoring level.}
Because of the recording drift, \neurhlG{}'s stat sheet is scaled by a declared
factor $m$ that multiplies regulation goal means (and so assists) and leaves
win probabilities unchanged. Before the season, $m_0 = L / M$, where $L$ is the
previous season's league regulation goals per team-game computed from the era
inputs the model already receives and $M$ is the engine's own mean projected
regulation goals over every game in the first 14 days, computed from roster
fallback lineups. In season,
\begin{equation}
m = \frac{A + kL}{P + kM}, \qquad k = 300 \text{ team-games},
\end{equation}
where $A$ and $P$ are the actual and unscaled predicted regulation goals over
completed 2026--27 games that have a primary forecast, so the factor moves
from $m_0$ toward the observed ratio as games accumulate and a forecast only
ever uses games already finished. On rosters from the day before the deadline
the factor was 0.922; it is fixed on the final rosters before the first game.

\paragraph{What the season can show.}
A single season cannot confirm effects of the size reported in
Section~\ref{sec:results}. Its purpose is prospective accountability: every
forecast is public before the event it forecasts, the scoring rules are fixed,
and the results will be published whatever they show. We will report them in
a revision of this paper after the regular season ends on 10 April 2027.
